% !PRINTED 632
\weberpage{632}
%<<CONT
\gesperrt{Grösse des Körpers $\Omega(\varepsilon)$ nur auf eine Weise in die Normalform
\[
\Phi(\varepsilon) = c_0 + c_1\varepsilon + c_2\varepsilon^2 + \cdots + c_{n-2}\varepsilon^{n-2}
\]
gebracht werden, worin die Co\"efficienten $c$ in $\Omega$ enthalten sind.}

Wir können in diesem Satze für $\Omega$ den Körper nehmen, der aus $R$ durch Adjunction von $n$ unabhängigen Variablen $x_0, x_1, x_2 \ldots x_{n-1}$ entsteht; denn wäre $X$ in diesem Körper reducibel, also
\[
X = X_1 X_2,
\]
wo $X_1, X_2$ ganze Functionen von $\varepsilon$ und rationale Functionen der $x$ sind, so könnte man für die Variablen $x_0, x_1 \ldots x_{n-1}$ solche rationale Zahlen setzen, dass $X_1, X_2$, ohne ihren Grad in Bezug auf $\varepsilon$ zu ändern, in Functionen in $R$ übergehen (\S.\,143), und dieses widerspricht der Irreducibilität von $X$ in $R$.

Dieser Körper $\Omega$ hat verschiedene Divisoren, in denen $X$ gleichfalls irreducibel ist, und die also alle im Theorem 1. für $\Omega$ genommen werden können. Solche Divisoren erhält man, wenn man irgend eine \gesperrt{Permutationsgruppe} $P$ der $n$ Ziffern $0, 1, 2 \ldots n-1$ festsetzt und den Inbegriff aller rationalen Functionen der $x_0, x_1 \ldots x_{n-1}$ in $R$ betrachtet, die die Permutationen dieser Gruppe \gesperrt{gestatten}.

Der Inbegriff aller dieser Functionen bildet offenbar einen Körper, und so bekommt man zu jeder Permutationsgruppe der $n$ Ziffern einen bestimmten zugehörigen Körper $\Omega$. Man kann die Grössen eines solchen Körpers rational darstellen durch die \gesperrt{symmetrischen Functionen der $x$ und durch eine zu der Gruppe $P$ gehörige Function}.

Wir nennen den so bestimmten Körper $\Omega$ \gesperrt{den zu der Gruppe $P$ gehörigen Körper}. Jeder solche Körper kann in 1. die Stelle von $\Omega$ vertreten. (Vgl. \S.\,152.)

Wir können also den zweiten Satz aufstellen:

\gesperrt{2. Wenn eine rationale Function $\Phi(x_0, x_1 \ldots x_{n-1}, \varepsilon)$ des Körpers $\Omega(\varepsilon)$ ungeändert bleibt, wenn auf die Indices der $x$ die Permutationen einer Gruppe $P$ angewandt werden, so gestatten auch die Coëfficienten $c_0, c_1, c_2 \ldots c_{n-2}$ der Normalform von $\Phi$ die Permutationen von $P$.}

Ist $\Phi(\varepsilon)$ irgend ein Element des Körpers $\Omega(\varepsilon)$, so erhält man die conjugirten Grössen, wenn man für $\varepsilon$ irgend eine andere
%CONT>>
